% Generated from locale/tr/content/computability/recursive-functions/other-recursions.tex; do not hand-edit.


\olfileid[tr]{cmp}{rec}{ore}
\olsection{Başka Özyineleme Biçimleri}

Çiftleme ve dizileme kullanarak ilkel özyinelemenin daha alışılmadık ve
yararlı biçimlerini gerekçelendirebiliriz. Örneğin aşağıdaki tanımda olduğu
gibi iki fonksiyonu aynı anda tanımlamak çoğu kez yararlıdır:
\begin{align*}
  h_0(\vec x,0) & = f_0(\vec x)\\
  h_1(\vec x,0) & = f_1(\vec x)\\
  h_0(\vec x,y+1) & =
    g_0(\vec x,y,h_0(\vec x,y),h_1(\vec x,y))\\
  h_1(\vec x,y+1) & =
    g_1(\vec x,y,h_0(\vec x,y),h_1(\vec x,y)).
\end{align*}
Bu, \emph{eşzamanlı özyineleme} örneğidir. Fonksiyon tanımlamanın bir başka
yararlı yolu da aşağıdaki gibi $h(\vec x,y+1)$ değerini
$h(\vec x,0),\dots,h(\vec x,y)$ değerlerinin \emph{tümü} cinsinden vermektir:
\begin{align*}
  h(\vec x,0) & = f(\vec x)\\
  h(\vec x,y+1) & =
    g(\vec x,y,\tuple{h(\vec x,0),\dots,h(\vec x,y)}).
\end{align*}
Aşağıdaki şema bu düşünceyi daha kısa biçimde anlatır:
\[
  h(\vec x,y)=
  g(\vec x,y,\tuple{h(\vec x,0),\dots,h(\vec x,y-1)}),
\]
burada $y=0$ iken $g$'nin son argümanı yalnızca boş dizidir. Her iki
anlatımda da düşünce, ``ardıl adımı'' hesaplanırken $h$ fonksiyonunun o ana
kadar hesaplanmış değerlerin bütün dizisini kullanabilmesidir. Buna
\emph{önceki değerler üzerinden özyineleme} denir. Özel bir örnekte bu,
aşağıdaki türden bir tanımı gerekçelendirmek için kullanılabilir:
\begin{align*}
  h(\vec x,y) & = \begin{cases}
    g(\vec x,y,h(\vec x,k(\vec x,y))) & \text{$k(\vec x,y)<y$ ise}\\
    f(\vec x) & \text{aksi durumda.}
  \end{cases}
\end{align*}
Başka bir deyişle $h$'nin $y$ noktasındaki değeri, $k$ tarafından verilen
\emph{herhangi bir} önceki noktadaki değeri cinsinden hesaplanabilir.

\begin{prob}
  Kalan fonksiyonu $r(x,y)$'yi önceki değerler üzerinden özyinelemeyle
  tanımlayın. ($x,y$ doğal sayılar ve $y>0$ ise $r(x,y)$, bazı $z$ için
  $x=z\times y+r(x,y)$ olacak biçimde $y$'den küçük sayıdır. Belirlilik için
  $y=0$ iken $r(x,0)=0$ diyelim.)
\end{prob}

Bu fonksiyonların sıradan ilkel özyineleme kullanılarak nasıl elde
edilebileceğini düşünmelisiniz. İlkel özyinelemenin son bir biçimi, süreç
boyunca \emph{parametrelerin} (yan değerlerin) değiştirilmesine izin verdiği
için daha esnektir:
\begin{align*}
  h(\vec x,0) & = f(\vec x)\\
  h(\vec x,y+1) & = g(\vec x,y,h(k(\vec x),y)).
\end{align*}
Bu biçim de sıradan ilkel özyinelemeyle benzetilebilir. Bunu yapmak kolay
değildir; ipucu olarak hesabı elle geriye doğru açmayı deneyin.

